\section{Limitations}
\label{sec:limits}

\subsection{Evaluation Scope}
\label{sec:stat_scope}
Each flight regime is evaluated on a \emph{single} representative
recording, so the headline numbers carry no between-flight error bars.  The
held-out flights (\autoref{sec:heldout}) and the cross-platform validation
(\autoref{sec:baselines}) broaden the evidence, but they do not substitute
for repeated trials.

The evaluation covers indoor flights on two platforms using the same
radar chip.  The main laboratory recordings last 18--26\,s.
Long-duration operation, different radar hardware, and performance
in smoke or fog remain untested.  The public flights use VINS pseudo
ground truth, and the backflip recording contains \ac{mocap} tracking
artifacts (\autoref{sec:setup}).

\subsection{Heading Sensor Requirement}
\label{sec:yaw}
Yaw is unobservable from Doppler$+$\ac{imu} (\autoref{sec:models}).
Here, reference yaw serves as an idealized magnetometer surrogate.
The evaluation therefore does not establish performance with a real
heading sensor, including its bias and time-varying magnetic
disturbances.  The weight is set at what a real sensor would justify:
a flight magnetometer resolves heading to a few degrees~\cite{kok2016magnetometer} with strongly
time-correlated errors, which under the correlation-time protocol of
\autoref{sec:char_budget} gives a weight of order $1$--$10$; we deploy
$\lambda_\psi = 0.6$, the conservative end of that range.  With
$\lambda_\psi = 0$ (no heading information after the start gauge) yaw
drifts within the flight: the orientation error grows to $10$ and
$9\degree$ on the two racing flights, with $0.5$--$0.6$\,m/s of
world-frame velocity error.  A ten times stronger weight reduces
orientation error by $0.03$--$0.2\degree$ on racing and $2.8\degree$
on the flip flight, but increases drift by up to $0.2$\,pp.
We retain the weaker weight to limit reliance on reference heading,
accepting the orientation trade-off.

\subsection{Extreme Dynamics: Backflips}
\label{sec:backflips}
The backflip flight circles in the horizontal plane while looping
continuously in the vertical plane at body rates around
$10$\,rad/s.  \autoref{sec:char} characterizes the
radar errors in this regime: 38\% of its returns alias, the noise floor of
the per-return law is four times the deployed nominal value, the
frame-shared error and the off-boresight excess the law does not carry
are at their largest.  One possible contributor is intra-frame
motion.  During a flip the scene sweeps some $17\degree$ across a
frame's integration time, and without
per-return timestamps the returns of one frame cannot be placed on the
trajectory individually.

The configuration also contributes to the backflip errors.
Robust accelerometer weighting improves orientation at the cost of
drift, while removing smoothness priors improves backflip accuracy
but worsens fast-racing drift (\autoref{sec:ablations}).  We retain
the shared configuration for its racing performance and report the
backflip degradation under the same settings.

\subsection{Sustained Flight Beyond the Doppler Ambiguity}
\label{sec:alias_limit}
Every aliased return has to be assigned a wrap. This is based on the last radar frame's consensus velocity, integrated forward with the
\ac{imu} (\autoref{sec:robust}). While only a few returns are aliased
the consensus rests on the unambiguous majority, but when most of the
frame is aliased it is built from the estimated unwraps themselves, and
an error can enter the consensus. This then displaces the next frame's
prediction by one wrap, $2v_{\max} = 6.3$\,m/s, and enters the cost
quadratically, so down-weighting does not help and a frame agreeing on
the same wrong wrap still passes the prefilter. The held-out circular
flight at 5\,m/s is this regime (\autoref{sec:heldout}).

\subsection{Vertical Channel}
\label{sec:vertical}
On fast racing the vertical channel carries most of the drift
(vertical \ac{rmse} $0.51$\,m against $0.27$\,m horizontally,
\autoref{fig:traj}).  On slow racing the corresponding errors are
$0.09$ and $0.11$\,m.  Sweeping the locked radar pitch from $25.5$ to
$29.5\degree$ moves the vertical drift monotonically in opposite
directions on the two racing flights.  This sweep does not identify
a common correction, so we keep the measured pitch.  The prefiltered
racing returns also show no constant elevation bias
(\autoref{sec:char_outliers}).

The fast-racing vertical error remains $0.51$\,m with or without
position-prior inflation.
The source of the remaining vertical drift is unresolved.

\subsection{Fixed Parameters and Motion Priors}
\label{sec:fixed_constants}
Two constants of the radar noise model are flight properties held
fixed at nominal-regime values: the noise floor $\sigma_0$ of the
per-return law (measured $0.049$--$0.327$\,m/s across regimes, deployed
at the racing geometric mean) and the frame-shared term $\sigma_c$
(measured $0.83$--$1.55$, deployed at $1.27$).  Both, like every
constant of the model, are properties of this radar and this
\ac{imu}: a different sensor requires the same measurements, and
\autoref{sec:char} is that procedure.  Online estimation
of these terms is a possible extension, not implemented here:
$\sigma_0$ is a within-frame statistic and can be measured in flight
without a reference, whereas $\sigma_c$ is a frame-mean statistic whose
in-flight estimate would have to separate the shared measurement error
from trajectory error.

The fixed motion priors, spline grids and window length
impose a trade-off between
racing and backflip performance (\autoref{sec:ablations}) and also
constrain the accelerometer bias estimates.  A motion prior derived
from measured motion statistics is one possible alternative to the
fixed smoothness penalties, but remains untested.

\subsection{Marginalization and Covariance}
\label{sec:marg_limits}
The marginalization prior is re-centered at each warm start with
its gradient set to zero, retaining only its curvature.  This does
not preserve the full Gaussian marginal.  Among factors touching
leaving states, those extending beyond the retained boundary are
dropped.  Angular-travel inflation approximates the effect of stale
linearizations (\autoref{sec:sw}).
The \ac{nees} results show conservative covariances on racing, but
do not establish consistency across regimes (\autoref{sec:nees}).

\subsection{Knot Density}
\label{sec:negative}
Finer orientation knots do not improve the tested backflip estimates.
At the deployed 16\,ms the orientation spline represents the \ac{mocap}
attitude through the flips to $1.3\degree$, the same at 8 and 4\,ms.
In the coarser-grid sweep, backflip orientation error is
$11.9$--$12.8\degree$ from 8 to 64\,ms and increases to $14.3\degree$
at 128\,ms, with $\lambda_\alpha$ held fixed across the sweep.
These results suggest that the orientation grid could be coarsened
where compute is tight.